\chapter{El álgebra de convolución de un grupo}\label{cap:v4-12}
La medida de Haar permite repetir la definición de convolución sin coordenadas. La ausencia de conmutatividad distingue el orden de los factores; la posible diferencia entre medidas izquierda y derecha introduce la función modular. Mantendremos visibles ambas cuestiones en las pruebas y en los ejercicios.

\section{Convolución e involución}\label{lectura:v4-22}
\index{convolucion@convolución!en grupos}\index{algebra@álgebra!L1 de un grupo}
Sea $G$ localmente compacto y Hausdorff, con medida de Haar izquierda $\mu$. Escribimos $dx=d\mu(x)$, y fijamos $\int h(xa)dx=\Delta(a)^{-1}\int h(x)dx$. Las clases $L^p$ se toman respecto de la completación de Haar. Para $f,g\in L^1(G)$ ponemos
\begin{equation}\label{v4-l24-d:algebra}
 (f*g)(x)=\int_Gf(y)g(y^{-1}x)dy,\qquad
 f^*(x)=\Delta(x^{-1})\overline{f(x^{-1})}.
\end{equation}
Recordamos la fórmula de inversión de Haar:
\begin{equation}\label{v4-l24-eq:inv}
 \int_Gh(x^{-1})\Delta(x^{-1})dx=\int_Gh(x)dx.
\end{equation}
En un grupo abeliano $\Delta=1$. En $\R^d$ recuperamos la convolución aditiva; en $\T$ usamos Haar de masa uno; en un grupo discreto con medida de conteo obtenemos $(a*b)(x)=\sum_y a(y)b(y^{-1}x)$. La suma puede tener índices no numerables, pero una función de $\ell^1$ tiene soporte numerable: para cada $n$ solo hay finitos términos de módulo al menos $1/n$, y la unión de esos conjuntos contiene todos los términos no nulos.

\begin{teorema}\label{v4:conv-grupo}\label{v4-l24-p:banach}
La convolución existe casi en todas partes y hace de $L^1(G)$ un álgebra de Banach con involución isométrica. En particular,
\[
 \norm{f*g}_1\leq\norm f_1\norm g_1,\quad
 \norm{f^*}_1=\norm f_1,\quad
 (f*g)^*=g^**f^*.
\]
\end{teorema}
\begin{proof}
Primero precisemos la cuestión de medida para grupos que no son $\sigma$-compactos. Cada función de $L^1$ tiene, salvo un nulo, soporte contenido en una unión numerable de compactos: se aproxima en $L^1$ por funciones de $C_c(G)$ y se elige una subsucesión con errores sumables. Una vecindad relativamente compacta y esa colección numerable de compactos generan un subgrupo abierto $\sigma$-compacto. En él la medida es $\sigma$-finita y se pueden aplicar Tonelli y Fubini ordinarios. Fuera de ese subgrupo las funciones elegidas y su convolución se anulan. Esto no afirma que el soporte topológico de todo representante sea $\sigma$-compacto.

Para representantes borelianos, $(y,x)\mapsto f(y)g(y^{-1}x)$ es medible. La transformación $(y,x)\mapsto(y,y^{-1}x)$ preserva la medida producto en el subgrupo anterior: integrar primero en $x$ y usar invariancia izquierda prueba la identidad para indicadoras y después para no negativas. Por ello sus preimágenes de nulos son nulas, y el argumento se extiende a representantes de la medida completada. Tonelli da
\[
 \int_G\int_G|f(y)g(y^{-1}x)|dydx
 =\int_G|f(y)|\int_G|g(z)|dzdy=\norm f_1\norm g_1.
\]
Fubini y la desigualdad triangular producen existencia, integrabilidad y la cota de norma. También prueban independencia de representantes y $\int f*g=(\int f)(\int g)$.

Para asociatividad, la integral triple absoluta es $\norm f_1\norm g_1\norm h_1$. El cambio $y=zt$ y Fubini dan
\[
 ((f*g)*h)(x)=\int\int f(z)g(z^{-1}y)h(y^{-1}x)dzdy
 =\int f(z)\int g(t)h(t^{-1}z^{-1}x)dtdz.
\]
La última expresión es $f*(g*h)(x)$ casi en todas partes. La bilinealidad es inmediata y la completitud es la de $L^1$.

La fórmula~\eqref{v4-l24-eq:inv} aplicada a $|f|$ prueba $\norm{f^*}_1=\norm f_1$. La antilinealidad y $f^{**}=f$ se obtienen de la conjugación y $\Delta(x^{-1})\Delta(x)=1$. Finalmente,
\begin{align*}
 (g^**f^*)(x)
 &=\int\Delta(y^{-1})\overline{g(y^{-1})}\Delta(x^{-1}y)\overline{f(x^{-1}y)}dy\\
 &=\Delta(x^{-1})\int\overline{g(y^{-1})}\,\overline{f(x^{-1}y)}dy\\
 &=\Delta(x^{-1})\int\overline{f(z)}\,\overline{g(z^{-1}x^{-1})}dz=(f*g)^*(x),
\end{align*}
donde en la última línea usamos $y=xz$. Las identidades se interpretan casi en todas partes.
\end{proof}

\section{Traslaciones y aproximaciones de la identidad}
\index{identidad!aproximada en un grupo}\index{red!de Dirac}
Definimos $L_af(x)=f(a^{-1}x)$ y $R_af(x)=f(xa)$. Un cambio de variable da
\[
 L_a(f*g)=(L_af)*g,\qquad R_a(f*g)=f*(R_ag),\qquad
 \norm{L_af}_1=\norm f_1,\quad\norm{R_af}_1=\Delta(a)^{-1}\norm f_1.
\]
\begin{lema}\label{v4-l24-l:translation}
Para cada $f\in L^1(G)$, $a\mapsto L_af$ es continua en norma. También lo es $a\mapsto\widetilde R_af$, donde $\widetilde R_af(x)=\Delta(a)^{-1}f(xa^{-1})$; estos últimos operadores son isometrías de $L^1$.
\end{lema}
\begin{proof}
Para $f\in C_c$ y $a$ en una vecindad compacta de $e$, los soportes de $f$ y de sus traslaciones quedan dentro de un compacto fijo. La continuidad conjunta, seguida de un recubrimiento finito de ese compacto, muestra que las diferencias tienden uniformemente a cero. Multiplicar la cota uniforme por la medida finita del compacto da continuidad en $L^1$. Para $f$ general, aproximamos por $g\in C_c$ y usamos
\[
 \norm{L_af-f}_1\leq2\norm{f-g}_1+\norm{L_ag-g}_1.
\]
En un punto distinto de $e$ se usa la ley de las traslaciones. Para $\widetilde R_a$, el mismo argumento usa continuidad de $\Delta$ y
$\norm{\widetilde R_af}_1=\Delta(a)^{-1}\Delta(a)\norm f_1=\norm f_1$.
\end{proof}
\begin{definicion}
Una red de Dirac es una red $(k_\alpha)\subset L^1(G)$ con $k_\alpha\geq0$, $\int k_\alpha=1$ y, para \emph{cada} vecindad $U$ de $e$,
\[
 \int_{G\setminus U}k_\alpha(x)dx\longrightarrow0.
\]
La condición de la lectura 22 que permite elegir alguna vecindad $U$ para cada $\eps$ es insuficiente: tomar $U=G$ la haría vacía de contenido.
\end{definicion}
\begin{teorema}\label{v4:dirac-grupo}
Para una red de Dirac y $f\in L^1(G)$, $k_\alpha*f\to f$ y $f*k_\alpha\to f$ en $L^1$.
\end{teorema}
\begin{proof}
Por Tonelli,
\[
 \norm{k_\alpha*f-f}_1\leq\int k_\alpha(y)\norm{L_yf-f}_1dy.
\]
La norma de la diferencia es pequeña en una vecindad de $e$ y está acotada por $2\norm f_1$ en todo $G$. La concentración de masa da el límite. Para el otro lado, el cambio $z=y^{-1}x$ y la fórmula de inversión dan
\[
 f*k_\alpha(x)=\int k_\alpha(z)\Delta(z)^{-1}f(xz^{-1})dz
 =\int k_\alpha(z)\widetilde R_zf(x)dz.
\]
La misma estimación con $\widetilde R_z$ prueba el segundo límite. El factor modular no puede omitirse en este paso.
\end{proof}
\begin{lema}\label{v4-l24-l:approx}
Existe una red $(e_U)\subset C_c(G)$, indexada por vecindades simétricas relativamente compactas de $e$ ordenadas por inclusión inversa, tal que
\[
 e_U\geq0,\quad\supp e_U\subset U,\quad e_U^*=e_U,
 \quad\int e_U=\norm{e_U}_1=1.
\]
Es una identidad aproximada por ambos lados.
\end{lema}
\begin{proof}
La local compacidad y la forma local de Urysohn proporcionan $h\in C_c(G)$ no negativa, no nula y apoyada en $U$. Su integral $a$ es positiva porque Haar da medida positiva a todo abierto no vacío. Sea $e_U=(h+h^*)/(2a)$. La involución preserva positividad e integral por~\eqref{v4-l24-eq:inv}, y $\supp h^*\subset U^{-1}=U$. Así se verifican las propiedades. Si $U$ se hace más pequeña que una vecindad fijada, el soporte queda dentro de ella, de modo que es una red de Dirac. El teorema anterior concluye.
Además se obtiene la cota útil
\begin{equation}\label{v4-l24-eq:approxbound}
 \norm{e_U*f-f}_1\leq\sup_{y\in U}\norm{L_yf-f}_1.
\end{equation}
\end{proof}
\begin{teorema}
El álgebra $L^1(G)$ es conmutativa si y solo si el grupo $G$ es abeliano.
\end{teorema}
\begin{proof}
Si $G$ es abeliano, $\Delta=1$ y el cambio $z=y^{-1}x$ transforma $f*g$ en $g*f$. Recíprocamente, sea conmutativa el álgebra. Para cada $a\in G$, los operadores de convolución por $L_ae_U$ son contracciones en $L^1$ y convergen fuertemente a $L_a$, porque $(L_ae_U)*f=L_a(e_U*f)\to L_af$. Todos esos operadores conmutan por hipótesis. Pasar a límites fuertes, primero en uno de los índices y después en el otro, usando sus cotas uniformes, da $L_aL_b=L_bL_a$, es decir, $L_{ab}=L_{ba}$.

La acción izquierda sobre $L^1$ es fiel: si $z\ne w$, podemos escoger $f\in C_c$ con $f(e)=1$ y $f(w^{-1}z)=0$ por separación local. En el punto $z$, $L_zf=1$ y $L_wf=0$; por continuidad difieren en un abierto de medida positiva. Por tanto $L_z\ne L_w$. Aplicar esto a $z=ab,w=ba$ prueba que $ab=ba$.
\end{proof}

\section{Ejercicios: modularidad y convolución}
\begin{ejercicio}[Lectura 22, ejercicio 1: Young con factor modular]\label{ej:v4-l22-1}
Sean $1/p+1/q=1+1/r$, con exponentes en $[1,\infty]$. Probar la versión válida en un grupo general:
\[
 \norm{f*g}_r\leq\norm f_p\norm{\Delta^{-1/p'}g}_q.
\]
En un grupo unimodular se reduce a la desigualdad usual. Explicar por qué la versión sin peso no vale siempre.
\end{ejercicio}
\begin{solucion}
Pongamos $h=\Delta^{-1/p'}g$. Para $r<\infty$ escribimos, con $z=y^{-1}x$,
\[
 |f(y)g(z)|=(|f(y)|^p|h(z)|^q)^{1/r}
 |f(y)|^{1-p/r}|h(z)|^{1-q/r}\Delta(z)^{1/p'}.
\]
Aplicamos Hölder con recíprocos $1/r$, $1/p-1/r$ y $1/q-1/r=1/p'$. Para el tercer factor, el cambio $y=xz^{-1}$ da
\[
 \int |h(y^{-1}x)|^q\Delta(y^{-1}x)dy=\int|h(z)|^qdz.
\]
Así
\[
 |f*g(x)|^r\leq\norm f_p^{r-p}\norm h_q^{r-q}(|f|^p*|h|^q)(x).
\]
El teorema~\ref{v4:conv-grupo}, seguido de integración, da la cota. Los factores de potencia cero se omiten. El caso $p=1$ se obtiene directamente por Jensen para $|f|/\norm f_1$; el caso $r=\infty$ es Hölder con el mismo cambio de variable y $q=p'$. De este modo se cubren los extremos sin dividir entre exponentes nulos.

El peso es real, no un artificio de la prueba. Para $1<p<\infty$ y $a$ con $\Delta(a)<1$, sean $g_U=L_ae_U$, de norma $L^1$ uno. Para $f\in C_c$ no nula, el cambio de variable anterior y continuidad de traslaciones muestran
\[
 f*g_U\longrightarrow\Delta(a)^{-1}f(\cdot a^{-1})\quad\text{en }L^p.
\]
Su norma límite es $\Delta(a)^{-1/p'}\norm f_p>\norm f_p$. La desigualdad sin peso, con $q=1,r=p$, fallaría para $U$ suficientemente pequeña. La convergencia se verifica expresando la integral como promedio de las funciones $\Delta(au)^{-1}f(\cdot u^{-1}a^{-1})$, $u\in U$, que convergen en $L^p$ a la función indicada; sobre una vecindad fija de $e$, las normas y factores modulares están uniformemente acotados.
\end{solucion}
\begin{ejercicio}[Lectura 22, ejercicio 2]\label{ej:v4-l22-2}
Si $G$ es discreto, demostrar que $L^1(G)$ tiene unidad.
\end{ejercicio}
\begin{solucion}
Por invariancia, todos los puntos tienen una misma masa $c>0$; es finita porque los puntos son compactos y positiva porque son abiertos. La función $u=c^{-1}\mathbf1_{\{e\}}$ satisface $u*f=f$ directamente en la suma de convolución. También $f*u=f$, pues solo contribuye el término $y=x$, de masa $c$. Con medida de conteo $c=1$ y la unidad es $\delta_e$. No se confunde esta función con una masa puntual como medida en un grupo no discreto.
\end{solucion}
\begin{proposicion}[El recíproco anunciado en la lectura]\label{v4:luno-unitaria}
Si $L^1(G)$ tiene unidad, $G$ es discreto.
\end{proposicion}
\begin{proof}
Primero, si $\mu(\{e\})=c>0$, una vecindad de clausura compacta tiene a lo sumo $\mu(K)/c$ puntos: de lo contrario contendría subconjuntos finitos con medida arbitrariamente grande. Esa vecindad es finita y abierta; quitando sus otros puntos, cerrados por Hausdorff, obtenemos que $\{e\}$ es abierto. Así $G$ es discreto.

Supongamos, para obtener una contradicción, que $G$ no es discreto y que $u\in L^1(G)$ es unidad. Entonces $\mu(\{e\})=0$. Por regularidad exterior y continuidad absoluta de la integral de $|u|$, hay una vecindad abierta $V$ de $e$ con $\int_V|u|<1$. Elijamos $\phi\in\Cc(G)$ con $0\le\phi\le1$, $\phi(e)=1$ y soporte contenido en $V^{-1}$; la existencia de estos cortes se probó al preparar Haar. La identidad $u*\phi=\phi$ en $L^1$ vale en cada punto: ambos lados tienen representantes continuos, por continuidad de convolución de $L^1$ con una función continua de soporte compacto, y una función continua cero casi en todas partes es cero en todas partes porque todo abierto no vacío tiene medida positiva. Evaluando en $e$,
\[
1=\left|\int_Gu(y)\phi(y^{-1})dy\right|\le\int_V|u(y)|dy<1,
\]
contradicción. La continuidad utilizada también se obtiene directamente por convergencia uniforme de las traslaciones de $\phi$, dominando la diferencia por $|u|$ multiplicado por el error uniforme.
\end{proof}
\begin{ejercicio}[Lectura 22, ejercicio 3: control de soportes]\label{ej:v4-l22-3}
Para $f\in C(G)$, demostrar que $e_U*f$ y $f*e_U$ existen en todo punto y convergen puntualmente a $f$, usando la red con soportes $\supp e_U\subset U$. Mostrar que, para una red de Dirac arbitraria, no basta que cada núcleo tenga soporte compacto.
\end{ejercicio}
\begin{solucion}
En la primera convolución, la variable de integración está en $\supp e_U$; en la segunda, está en $x(\supp e_U)^{-1}$. Son compactos y $f$ es acotada en ellos, así que ambas integrales existen. Para $x$ fijo,
\[
 |e_U*f(x)-f(x)|\leq\sup_{y\in U}|f(y^{-1}x)-f(x)|\longrightarrow0.
\]
Para la otra convolución usamos
\[
 f*e_U(x)=\int e_U(z)\Delta(z)^{-1}f(xz^{-1})dz.
\]
La función de $z$ dentro del promedio es continua y vale $f(x)$ en $e$; el mismo argumento prueba su límite.

La concentración de masa sola no controla una función no acotada. En $\R$, sea $f(x)=e^{x^2}$, sea $b$ una función suave no negativa de integral uno apoyada en $[-1,1]$, y sea $h_n$ una identidad aproximada de soporte $[-1/n,1/n]$. Entonces $k_n=(1-1/n)h_n+n^{-1}b(\cdot-n)$ es una sucesión de Dirac, cada término de soporte compacto, pero
\[
 (k_n*f)(0)\geq n^{-1}e^{(n-1)^2}\longrightarrow\infty.
\]
Por tanto se corrige el enunciado usando soportes que se contraen, condición que cumple la red construida en este capítulo.
\end{solucion}
\begin{ejercicio}[Lectura 22, ejercicio 4]\label{ej:v4-l22-4}
Si $f\in L^1(G)$ y $g\in L^\infty(G)$, probar que $f*g$ está definida en todo punto, es acotada y satisface $\norm{L_a(f*g)-f*g}_\infty\to0$ cuando $a\to e$.
\end{ejercicio}
\begin{solucion}
Elijamos un representante acotado de $g$. Entonces la integral absoluta es a lo sumo $\norm f_1\norm g_\infty$ para cada $x$. Otro representante equivalente da la misma integral: inversión y traslaciones preservan nulos porque el factor modular es positivo. La identidad $L_a(f*g)=(L_af)*g$ da
\[
 \norm{L_a(f*g)-f*g}_\infty\leq\norm{L_af-f}_1\norm g_\infty\longrightarrow0.
\]
Esta es la continuidad uniforme por traslaciones izquierdas que se afirma; en particular implica continuidad ordinaria. En un grupo no abeliano conviene especificar la uniformidad utilizada en lugar de intercambiar sin explicación izquierda y derecha.
\end{solucion}
\begin{ejercicio}[Lectura 22, ejercicio 5]\label{ej:v4-l22-5}
Si $f,g\in C_c(G)$, demostrar $\supp(f*g)\subseteq\supp f\,\supp g$ y que $f*g\in C_c(G)$.
\end{ejercicio}
\begin{solucion}
Si $x$ no pertenece al producto de soportes, no puede ocurrir simultáneamente $y\in\supp f$ y $y^{-1}x\in\supp g$; el integrando es cero. El producto de soportes es compacto por continuidad de multiplicación. La convolución es continua por el ejercicio anterior, pues $g$ es acotada y $f\in L^1$. Así se anula fuera de un compacto y pertenece a $C_c$.
\end{solucion}

\paragraph{Los operadores que produce la convolución.}
La unidad puede faltar en $L^1(G)$, pero la red $(e_U)$ permite aproximar su acción. Esa red hará posible reconstruir operadores del grupo a partir de una representación del álgebra, que es el objetivo del capítulo siguiente.
